\documentclass[12pt,a4paper]{article}

% Compile with: tectonic Lab18_Moisei_Daniel_FAF241_Report.tex
\usepackage{fontspec}
\setmainfont{Liberation Serif}
\setmonofont{Liberation Mono}
\usepackage[left=2cm,right=1cm,top=2cm,bottom=2cm]{geometry}
\usepackage{setspace}
\usepackage{graphicx}
\usepackage[export]{adjustbox}
\usepackage{xcolor}
\usepackage{listings}
\usepackage{titlesec}
\usepackage{array}
\usepackage{tabularx}
\usepackage{caption}
\usepackage{float}
\usepackage{chngcntr}
\usepackage[hidelinks]{hyperref}

\graphicspath{{figures/}}
\counterwithin{figure}{section}
\definecolor{codebackground}{HTML}{F2F2F2}
\setlength{\parindent}{1.25cm}
\setlength{\parskip}{0pt}
\onehalfspacing
\setlength{\intextsep}{8pt plus 2pt minus 2pt}
\captionsetup[figure]{font=small,skip=4pt,justification=centering}

\titleformat{\section}
  {\normalfont\bfseries\centering\fontsize{13}{16}\selectfont}
  {\thesection}{0.8em}{\MakeUppercase}
\titleformat{name=\section,numberless}
  {\normalfont\bfseries\centering\fontsize{13}{16}\selectfont}
  {}{0pt}{\MakeUppercase}
\titleformat{\subsection}
  {\normalfont\bfseries\fontsize{12}{15}\selectfont}
  {\thesubsection}{0.7em}{}
\titlespacing*{\section}{0pt}{0pt}{12pt}
\titlespacing*{\subsection}{0pt}{10pt}{6pt}

\lstset{
  basicstyle=\ttfamily\small,
  backgroundcolor=\color{codebackground},
  frame=none,
  breaklines=true,
  columns=fullflexible,
  keepspaces=true,
  showstringspaces=false,
  keywordstyle=\ttfamily,
  commentstyle=\ttfamily,
  stringstyle=\ttfamily,
  xleftmargin=0.4cm,
  xrightmargin=0.2cm,
  aboveskip=6pt,
  belowskip=10pt
}
\newcommand{\taskfigure}[4][]{%
  \begin{figure}[H]
    \centering
    \includegraphics[scale=0.52,max width=\textwidth,max height=0.85\textheight]{figures/#2}
    \caption{#3}
    \label{fig:#4}
  \end{figure}
}
\newcommand{\inlinecmd}[1]{\texttt{\detokenize{#1}}}

\begin{document}
\sloppy

\begin{titlepage}
  \begin{center}
    \includegraphics[width=3.4cm]{figures/LOGO_UTM.jpg}\par
    \vspace{0.25cm}

    {\bfseries\fontsize{14}{17}\selectfont
      MINISTRY OF EDUCATION, CULTURE AND RESEARCH\par
      OF THE REPUBLIC OF MOLDOVA\par
      Technical University of Moldova\par
      Faculty of Computers, Informatics and Microelectronics\par
      Department of Software Engineering and Automation\par
    }

    \vspace{1.25cm}
    {\bfseries\fontsize{16}{19}\selectfont Daniel Moisei FAF-241\par}

    \vspace{1.35cm}
    {\bfseries\fontsize{34}{39}\selectfont Report\par}
    \vspace{0.2cm}
    {\itshape\fontsize{20}{24}\selectfont Laboratory Work No. 18\par}
    {\bfseries\itshape\fontsize{20}{24}\selectfont Operating Systems\par}

    \vspace{0.8cm}
    {\fontsize{18}{22}\selectfont Topic: A User-Level Thread Switch (xv6)\par}
  \end{center}

  \vfill
  \begin{flushright}
    {\fontsize{15}{18}\selectfont Checked by:}\par
    \textbf{Reițman Patricia}, \textit{university assistant}
  \end{flushright}

  \vfill
  \begin{center}
    {\bfseries\fontsize{14}{17}\selectfont Chișinău - 2026}
  \end{center}
\end{titlepage}

\pagenumbering{arabic}
\setcounter{page}{2}
\tableofcontents
\newpage

\section*{Introduction}
\addcontentsline{toc}{section}{Introduction}

The purpose of this laboratory work was to build a tiny thread library in user space, inside one xv6 program. The program \texttt{uthread} creates three threads and switches between them by saving and restoring registers, exactly like the kernel function \texttt{swtch}, but entirely in its own code. The kernel was not changed: for xv6, \texttt{uthread} is one ordinary process, and the three threads exist only inside it.

Three edits were needed: the thread library in \texttt{user/uthread.c}, the switch itself in assembly in \texttt{user/uthread\_switch.S}, and a rule in the \texttt{Makefile} that links the two together.

\newpage
\section{The Switch}

\subsection{The thread library}

Each thread has a saved context, its own stack of 8\,KB, and a state (\texttt{FREE}, \texttt{RUNNING} or \texttt{RUNNABLE}). The skeleton of the handout had three places to complete:

\begin{itemize}
  \item \texttt{struct context}: the 14 callee-saved registers \texttt{ra}, \texttt{sp} and \texttt{s0} to \texttt{s11}, each a \texttt{uint64}, in the order in which \texttt{thread\_switch} stores them.
  \item \texttt{thread\_create}: the context of a new thread gets \texttt{ra} set to the thread function and \texttt{sp} set to the top of the thread's stack, because the stack grows down. The first switch to the thread then ``returns'' into its function, on its own stack.
  \item \texttt{thread\_schedule}: the next thread is marked \texttt{RUNNING}, the old thread is remembered, the next thread becomes the current one, and \texttt{thread\_switch(\&old->ctx, \&next\_thread->ctx)} is called.
\end{itemize}

\texttt{thread\_b} and \texttt{thread\_c} are copies of \texttt{thread\_a} in which \texttt{a} was changed to \texttt{b} or \texttt{c}; each waits until the other two have started.

\lstinputlisting[language=C,caption={Implementation of user/uthread.c}]{../uthread.c}

\subsection{thread\_switch}

\texttt{thread\_switch(old, new)} has the same shape as the kernel's \texttt{swtch}: 14 \texttt{sd} instructions save \texttt{ra}, \texttt{sp} and \texttt{s0} to \texttt{s11} of the current thread at the offsets 0 to 104 of \texttt{old} (register \texttt{a0}), 14 \texttt{ld} instructions load the same registers of the next thread from \texttt{new} (register \texttt{a1}), and \texttt{ret} returns into the next thread.

\lstinputlisting[caption={Implementation of user/uthread\_switch.S}]{../uthread_switch.S}

\subsection{The Makefile}

The line \texttt{\$U/\_uthread\textbackslash} was added to \texttt{UPROGS}, and the rule below was added at the end of the \texttt{Makefile}, so that \texttt{uthread.o} is linked together with \texttt{uthread\_switch.o}:

\noindent\begin{minipage}{\linewidth}
\begin{lstlisting}
$U/_uthread: $U/uthread.o $U/uthread_switch.o $(ULIB) $U/user.ld
	$(LD) $(LDFLAGS) -T $U/user.ld -o $U/_uthread $U/uthread.o $U/uthread_switch.o $(ULIB)
	$(OBJDUMP) -S $U/_uthread > $U/uthread.asm
\end{lstlisting}
\end{minipage}

\newpage
\section{Run It}

\subsection{The output}

After rebuilding, \texttt{uthread} printed exactly the expected output. The beginning is shown in Figure~\ref{fig:start}: each thread prints its \texttt{started} line and yields until the other two have started, then the three take turns, \texttt{c}, \texttt{a}, \texttt{b}. The order is round robin, because every \texttt{thread\_yield} switches to the next runnable thread after the current one. \texttt{thread\_c} counts first because when it starts it finds \texttt{a} and \texttt{b} already started and goes straight into its loop.

\taskfigure{part2-uthread-start.png}{The beginning: the three threads start and take turns}{start}

The end is shown in Figure~\ref{fig:end}. Each thread counted from 0 to 99 and exited after 100 steps. When the last one had exited, no thread was runnable and \texttt{thread\_schedule} ended the program with \texttt{thread\_schedule: no runnable threads}.

\taskfigure{part2-uthread-end.png}{The end: every thread exits after 100}{end}

\subsection{Answer}

\texttt{thread\_switch} saves the 14 callee-saved registers of the running thread into \texttt{old}: \texttt{ra}, the return address, \texttt{sp}, the stack pointer, and \texttt{s0} to \texttt{s11}. Then it loads the same 14 registers of the next thread from \texttt{new} and executes \texttt{ret}. These registers are enough because \texttt{thread\_switch} is called like a normal C function: the compiler has already saved every caller-saved register (\texttt{t0}--\texttt{t6}, \texttt{a0}--\texttt{a7}) that the caller still needs, on the caller's own stack. Loading \texttt{sp} moves to the stack of the other thread, where its own saved values and local variables are.

\texttt{ra} makes execution continue in the other thread because \texttt{ret} means ``jump to the address in \texttt{ra}''. After the loads, \texttt{ra} holds the return address of the next thread, so \texttt{ret} does not return into the thread that called \texttt{thread\_switch}, but into the other one. For a thread that has already run, this is the point right after its own call to \texttt{thread\_switch} inside \texttt{thread\_schedule}. For a new thread it is the first instruction of its function, because \texttt{thread\_create} set \texttt{ctx.ra} to the function and \texttt{ctx.sp} to the top of the thread's stack.

\newpage
\section*{Conclusion}
\addcontentsline{toc}{section}{Conclusion}

This laboratory work showed that switching threads needs no help from the kernel: a thread is only a set of saved registers and its own stack. Saving the callee-saved registers of one thread and loading those of another, then executing \texttt{ret}, is enough to continue in a different thread, because the loaded \texttt{ra} decides where \texttt{ret} jumps and the loaded \texttt{sp} decides which stack is used. A new thread is started by preparing these two registers by hand, so that the first switch ``returns'' into the start of its function.

The three threads took turns exactly as expected, in round-robin order, each giving up the CPU voluntarily with \texttt{thread\_yield}. This is the same mechanism the xv6 kernel uses in \texttt{swtch} to switch between processes, with one difference: here the switch happens only when a thread calls \texttt{thread\_yield}, since there is no timer interrupt to take the CPU away from a thread that does not yield.

\end{document}
